\section{Canvas interactions}
\label{sec:canvas}

The centre of the window is a QGraphicsView-based canvas with
layered overlays. From bottom to top:

\begin{enumerate}
  \item Background image (reference or deformed).
  \item Mesh overlay (when enabled).
  \item Region of Interest outline.
  \item Field overlay (post-run, colour-mapped displacement or
        strain).
  \item Refine-brush overlay (cyan, frame 1 only).
\end{enumerate}

\subsection{Zoom and pan}

\begin{description}[leftmargin=1.2cm,style=nextline]
  \item[Mouse wheel] Zoom in / out at the cursor position.
  \item[Fit] Button at the top of the canvas toolbar --- zoom to fit
    the full image in the viewport.
  \item[100\%] Zoom to 1:1 pixel ratio.
  \item[+ / --] Buttons for discrete zoom steps.
  \item[Middle-click drag] Pan the view.
  \item[Left-click drag in empty area] Pan the view (if no drawing
    tool is active).
\end{description}

\subsection{Frame navigation}

The frame slider at the bottom of the canvas scrubs through the
sequence. The label to the left shows the current frame number and
total count.

A small play/pause button next to the slider auto-advances at a
configurable frame rate --- useful for reviewing displacement fields
after a run.

\subsection{Show on deformed frame}

This toggle lives in the FIELD section of the right sidebar (not
VISUALIZATION, because it controls \emph{where} the field plots,
i.e.\ geometry, rather than how it styles).

\begin{description}[leftmargin=1.2cm,style=nextline]
  \item[Checked (default)]
    The field overlay is plotted on the currently-selected deformed
    frame, with node positions warped by the cumulative
    displacement.
  \item[Unchecked]
    The field overlay is plotted on the reference frame (frame 1)
    with node positions un-warped. Useful for comparing an
    experiment back to the starting configuration.
\end{description}
